% =============================================================================
% Figure placeholders.
%
% Each block below carries the FINAL caption and the \label used in the text.
% Replace \placeholderbox{...} with \includegraphics{...} once the artwork
% exists. Generation recipes: paper/figures/README.md.
%
% The three Method figures (bridge architecture, vehicle, pipeline) live in
% sections/method.tex so they float near their citations. \placeholderbox is
% defined in main.tex.
% =============================================================================




\begin{figure}[t]
\centering
\placeholderbox{Six stacked panels sharing a \SI{30}{\second} time axis.}{5.2cm}
\caption{Example telemetry from one identification run. (a) $v_{x}$;
(b) $v_{y}$ and $\omega$; (c) commanded versus achieved road-wheel angle,
whose ratio measures \num{0.76}; (d) $a_{x}$, $a_{y}$; (e) per-axle lateral
force from the simulator's wheel telemetry against the value recovered by the
steady-state relation~\eqref{eq:ssforces}; (f) per-axle slip angle.}
\label{fig:telemetry}
\end{figure}

\begin{figure}[t]
\centering
\placeholderbox{$\mu_{\mathrm{reach}}$ versus rollout speed, one curve per
sweep amplitude, with the old clip band and the plant's measured
$\mu$.}{4.6cm}
\caption{Reachable friction of the synthetic rollout,
$\mu_{\mathrm{reach}}=v^{2}\delta_{\mathrm{sweep}}/(gL)$, for the full-scale
vehicle ($L=\SI{1.53}{\meter}$, solid) and a 1:10 platform
($L=\SI{0.32}{\meter}$, dashed). The shaded band is the rollout speed
inherited from the scaled platform; within it the full-scale rollout cannot
demand more than $\mu=0.43$, so the tire's peak factor is not identifiable
from it. The same configuration is comfortably above the peak at 1:10, which
is why the defect is scale-latent.}
\label{fig:mureach}
\end{figure}

\begin{figure}[t]
\centering
\placeholderbox{(a) Slip-angle histograms before/after; (b) $(\alpha,F_{y})$
scatter with fitted curve, before/after.}{4.4cm}
\caption{Excitation entering the Pacejka fit. (a) Distribution of the slip
angles present in the synthetic data before and after the rollout correction,
with the reference tire's peak-force slip angle marked. (b) The corresponding
$(\alpha,F_{y})$ scatter with the fitted curve; before the correction all data
lies on the initial linear ramp, where only the product $BCD$ is
identifiable.}
\label{fig:excitation}
\end{figure}

\begin{figure}[t]
\centering
\placeholderbox{$2\times4$ grid: front/rear $\times$ $B,C,D,E$ versus
co-identification iteration, two series.}{4.8cm}
\caption{Coefficient trajectories over the six co-identification iterations,
with the unbounded sweep and tracking regularisation target (which walks
$D$ from \num{1.01} to \num{1.87}) against the bounded sweep with a frozen
target (which settles in $[1.05,1.18]$). Dashed lines mark the box
constraints and the reference values of Table~\ref{tab:ref}.}
\label{fig:iterations}
\end{figure}

\begin{figure}[t]
\centering
\placeholderbox{Two panels (front, rear): $F_{y}(\alpha)$ for reference,
degenerate, and corrected fits over the measured scatter.}{4.4cm}
\caption{Identified tire models against the plant. Each panel shows the
reference fit of Table~\ref{tab:ref}, the degenerate result of the
untransplanted pipeline, and the corrected result over three consecutive
identifications (shaded band), with the raw per-wheel
$(\alpha,F_{y})$ measurements underneath.}
\label{fig:tirecurves}
\end{figure}

\begin{figure}[t]
\centering
\placeholderbox{Measured versus predicted $v_{y}$ and $\omega$ on held-out
data, with residual histograms inset.}{4.2cm}
\caption{One-step prediction on a held-out run using the corrected identified
model: $R^{2}=\num{0.97}$ for lateral velocity and $R^{2}=\num{0.99}$ for yaw
rate.}
\label{fig:onestep}
\end{figure}

\begin{figure}[t]
\centering
\placeholderbox{Pending: demand profile after regenerating the raceline for
the measured vehicle, on the same axes as Fig.~\ref{fig:raceline}(c).}{4.0cm}
\caption{Regenerated speed profile. \textbf{Not yet produced}: the
demand-versus-grip comparison for the \emph{current} trajectory is
Fig.~\ref{fig:raceline}(c); this figure is its counterpart after the raceline
is regenerated for the measured vehicle ($m=\SI{269.6}{\kilo\gram}$,
$\mu_{\mathrm{peak}}=1.0$), which is the prerequisite for the closed-loop
experiment of Fig.~\ref{fig:closedloop}.}
\label{fig:demand}
\end{figure}

\begin{figure}[t]
\centering
\placeholderbox{Grouped bars: $C_{f}$, $C_{r}$ across five model
sources.}{3.8cm}
\caption{Axle cornering stiffness $C=F_{z}BC^{\mathrm{MF}}D$ across model
sources: simulator ground truth, the degenerate fit, a fit referenced to the
steering command, a fit referenced to the achieved steering angle, and the
final corrected fit. The understeer condition $l_{f}C_{f}$ versus
$l_{r}C_{r}$ is annotated per group.}
\label{fig:stiffness}
\end{figure}

\begin{figure}[t]
\centering
\placeholderbox{Pending: cross-track error versus time, and driven path over
the raceline.}{3.8cm}
\caption{Closed-loop tracking under MPC with the identified model.
\textbf{Not yet produced}: this experiment is blocked on regenerating a
feasible raceline for the measured vehicle, see
Section~\ref{subsec:negative}.}
\label{fig:closedloop}
\end{figure}

\begin{figure*}[t]
\centering
\placeholderboxwide{Pending: three panels: (a) $\mu_{\mathrm{reach}}$ versus
rollout speed at full scale and 1:10; (b) round-trip identified $D$ versus
rollout speed under the baseline and the excitation-aware configuration; (c)
the $(\alpha_{f},F_{y,f})$ coverage each configuration hands to the
fit.}{5.0cm}
\caption{The baseline pipeline~\cite{ontrack2024} against this work, with the
residual network held at zero so that the rollout configuration is the only
variable. \textbf{Not yet produced}; generation recipe in
\texttt{figures/make\_baseline\_comparison.py}. (a) The friction the synthetic
rollout can demand, \eqref{eq:mureach}: the baseline's
\SIrange{2}{4}{\meter\per\second} clip reaches $\mu=0.43$ on the full-scale
vehicle but $\approx2.0$ on the 1:10 platform it was designed for, which is
why the defect is scale-latent. (b) A known tire rolled out and fitted back
across rollout speeds: inside the baseline clip the peak factor lands on its
box bound, above \SI{7}{\meter\per\second} it is recovered. (c) The
$(\alpha_{f},F_{y,f})$ data the Pacejka fit actually sees under each
configuration, against the true curve; the baseline configuration supplies
only the initial linear ramp, where solely the product $BCD$ is
identifiable.}
\label{fig:baseline}
\end{figure*}
